\begin{figure*}
\centering
\captionsetup[subfigure]{font=footnotesize,skip=0pt}

\begin{subfigure}{\linewidth}
\centering
\includegraphics[width=0.9\linewidth]{../binder/binder-2026-07-06-pcomplexity-gee.ipynb/binder/teeplots/2026-07-06-pcomplexity-gee/other=co-evo+viz=subplots+ext=.png}
\caption{fitness in co-evolved vs. self context, for original 100-stint replicates; markers indicate significance at $p<0.05,0.01,0.001$ (Mann-Whitney tests)}
\label{fig:general-supp:eco-vs-self}
\end{subfigure}

% https://github.com/mmore500/oee4/blob/790894365e523ecf6af77060c8b2bc76f9df5cc4/binder/2026-07-06-pcomplexity-gee.ipynb
\begin{subfigure}[t]{0.32\linewidth}
\centering
\includegraphics[width=\linewidth]{../binder/binder-2026-07-06-pcomplexity-gee.ipynb/binder/teeplots/2026-07-06-pcomplexity-gee/other=co-evo+viz=subplots+what=exclude+ext=.pdf}
\caption{GEE, excluding case study replicate 16005 (interaction coef. 0.1455)}
\label{fig:general-supp:exclude-co-evo}
\end{subfigure}
\hfill
\begin{subfigure}[t]{0.32\linewidth}
\centering
\includegraphics[width=\linewidth]{../binder/binder-2026-07-06-pcomplexity-gee.ipynb/binder/teeplots/2026-07-06-pcomplexity-gee/other=co-evo+viz=subplots+what=disagg+ext=.pdf}
\caption{GEE, disaggregated by stint within replicates (interaction coef. 0.0477; autoregressive param. 0.194)}
\label{fig:general-supp:disagg-co-evo}
\end{subfigure}
\hfill
\begin{subfigure}[t]{0.32\linewidth}
\centering
\includegraphics[width=\linewidth]{../binder/binder-2026-07-06-pcomplexity-gee.ipynb/binder/teeplots/2026-07-06-pcomplexity-gee/other=co-evo+viz=subplots+what=disagg-excl+ext=.pdf}
\caption{GEE, disaggregated by stint within replicates excluding case study replicate 16005 (interaction coef. 0.0464; autoregressive param. 0.2)}
\label{fig:general-supp:disagg-excl-co-evo}
\end{subfigure}

\begin{subfigure}{0.49\linewidth}
\centering
\includegraphics[width=\linewidth]{../binder/binder-2026-07-25-ecoselfcontext-pgcomplex-screen.ipynb/binder/teeplots/2026-07-25-ecoselfcontext-pgcomplex-screen/inner=box+iqr=shade+n=10+viz=subplots+ext=.pdf}
\caption{top 10 replicates by g/p complexity joint exceedance (shaded)}
\label{fig:general-supp:pgcomplex-ecoself-n10}
\end{subfigure}
\begin{subfigure}{0.49\linewidth}
\centering
\includegraphics[width=\linewidth]{../binder/binder-2026-07-25-ecoselfcontext-pgcomplex-screen.ipynb/binder/teeplots/2026-07-25-ecoselfcontext-pgcomplex-screen/inner=box+iqr=shade+n=30+viz=subplots+ext=.pdf}
\caption{top 30 replicates by g/p complexity joint exceedance (shaded)}
\label{fig:general-supp:pgcomplex-ecoself-n30}
\end{subfigure}

\begin{subfigure}{0.32\linewidth}
\centering
\includegraphics[width=\linewidth]{../binder/binder-2026-07-25-ecoselfcontext-pgcomplex-screen.ipynb/binder/teeplots/2026-07-25-ecoselfcontext-pgcomplex-screen/alt=nobb+inner=box+iqr=shade+n=10+viz=subplots+ext=.pdf}
\caption{top 10, baseline as no bkgd.}
\label{fig:general-supp-nobb:n10}
\end{subfigure}
\begin{subfigure}{0.32\linewidth}
\centering
\includegraphics[width=\linewidth]{../binder/binder-2026-07-25-ecoselfcontext-pgcomplex-screen.ipynb/binder/teeplots/2026-07-25-ecoselfcontext-pgcomplex-screen/alt=nobb+inner=box+iqr=shade+n=20+viz=subplots+ext=.pdf}
\caption{top 20, baseline as no bkgd.}
\label{fig:general-supp-nobb:n20}
\end{subfigure}
\begin{subfigure}{0.32\linewidth}
\centering
\includegraphics[width=\linewidth]{../binder/binder-2026-07-25-ecoselfcontext-pgcomplex-screen.ipynb/binder/teeplots/2026-07-25-ecoselfcontext-pgcomplex-screen/alt=nobb+inner=box+iqr=shade+n=30+viz=subplots+ext=.pdf}
\caption{top 30, baseline as no bkgd.}
\label{fig:general-supp-nobb:n30}
\end{subfigure}

% https://github.com/mmore500/oee4/blob/f996435b781cad8f3ff051a6140c71a0ca165966/binder/2026-07-25-ecoselfcontext-pgcomplex-screen.ipynb

\caption{%
\textbf{Generality of Relationship Between Biogenic Niche Construction and Complexity.}
\footnotesize
Biogenic niche construction effects were detected in 11/40 replicates of original 100-stint trials (panel \ref{fig:general-supp:eco-vs-self}, case study replicate leftmost; see also Supplementary Figure \ref{fig:external-bb-all}).
Across GEE test designs, phenotype complexity consistently drove self-context fitness increases compared to co-evolved context fitness (panels \labelcref{fig:general-supp:exclude-co-evo,fig:general-supp:disagg-co-evo,fig:general-supp:disagg-excl-co-evo}).
95\% CI for interaction coefficients are 0.1455 (0.049, 0.242), 0.0477 (0.010, 0.086), and 0.0464 ($9.19\times10^{-5}$, 0.093), respectively; for main text panel \ref{fig:general:gee-base}, 0.1101 (0.048, 0.172).
Reported tests were performed at stints 9, 19, 29, 39, 49, 59, 69, 79, 89, 99 --- replicating exploratory tests performed over stints 10, 20, 30, 40, 50, 60, 70, 80, 90, 100.
Phenotype complexity was chosen over genotype complexity, due to better comparability between strains and over time (fixed number of trials per genome, across possible I/O types).
GEE test takes form \texttt{Q("Fitness") \textasciitilde{} C("Fitness Kind") * Q("Phenotype Complexity")} with binomial family configuration.
Bands denote 95\% confidence interval.
Association between high-complexity outcomes in additional 20-stint replicates and biogenic niche construction (fitness advantage in self vs. co-evolved context) was consistently detected at 2.5\% and 7.5\% thresholds of joint exceedance of genotype and phenotype complexity (Fisher's exact test, $p=0.027$, $p=0.005$; 9/1 vs. 212/78, 26/7 vs. 195/172; panels \labelcref{fig:general-supp:pgcomplex-ecoself-n10,fig:general-supp:pgcomplex-ecoself-n30}).
Panels \labelcref{fig:general-supp-nobb:n10,fig:general-supp-nobb:n20,fig:general-supp-nobb:n30} show that the association also appears when comparing self-context fitness to no-background baseline fitness, rather than baseline fitness in co-evolved context (Fisher's exact test, $p = 0.182$, $p = 0.006$, $p = 0.011$; 9/1 vs. 266/124, 19/1 vs. 256/124, 29/4 vs. 246/121).
Note low statistical sensitivity at $n=10$ (panel \ref{fig:general-supp-nobb:n10}) due to a higher fraction of overall replicates exhibiting higher self-context fitness than no-background fitness.
Supplement to Figure \ref{fig:general}.
}
\label{fig:general-supp}

\end{figure*}

% \begin{subfigure}{0.49\linewidth}
% \centering
% \includegraphics[width=\linewidth]{../binder/binder-2026-07-06-pcomplexity-gee.ipynb/binder/teeplots/2026-07-06-pcomplexity-gee/other=none+viz=subplots+what=exclude+ext=.pdf}
% \caption{excluding case-study replicate 16005, self vs.\ no background}
% \label{fig:general-supp:exclude-none}
% \end{subfigure}

% \begin{subfigure}{0.49\linewidth}
% \centering
% \includegraphics[width=\linewidth]{../binder/binder-2026-07-06-pcomplexity-gee.ipynb/binder/teeplots/2026-07-06-pcomplexity-gee/other=none+viz=subplots+what=disagg+ext=.pdf}
% \caption{stints disaggregated, self vs.\ no background}
% \label{fig:general-supp:disagg-none}
% \end{subfigure}

% \begin{subfigure}{0.49\linewidth}
% \centering
% \includegraphics[width=\linewidth]{../binder/binder-2026-07-06-pcomplexity-gee.ipynb/binder/teeplots/2026-07-06-pcomplexity-gee/other=none+viz=subplots+what=disagg-excl+ext=.pdf}
% \caption{stints disaggregated, excluding replicate 16005, self vs.\ no background}
% \label{fig:general-supp:disagg-excl-none}
% \end{subfigure}

% \begin{subfigure}{0.49\linewidth}
% \centering
% \includegraphics[width=\linewidth]{../binder/binder-2026-07-06-pcomplexity-gee.ipynb/binder/teeplots/2026-07-06-pcomplexity-gee/other=none+viz=subplots+what=base+ext=.pdf}
% \caption{base condition, self vs.\ no background}
% \label{fig:general-supp:base-none}
% \end{subfigure}
% \hfill

% \begin{subfigure}{\linewidth}
% \centering
% \includegraphics[width=0.9\linewidth]{../binder/binder-2026-07-06-pcomplexity-gee.ipynb/binder/teeplots/2026-07-06-pcomplexity-gee/other=none+viz=subplots+ext=.png}
% \caption{self vs.\ no-background fitness prevalence, compared across independent replicates}
% \label{fig:general-supp:overview-none}
% \end{subfigure}
